\documentclass{article}
\usepackage{amsmath,amssymb,xurl}
\usepackage[hidelinks]{hyperref}
% Shared commands for the notes in docs/paper-gaps/.
% Every note loads this file via % Shared commands for the notes in docs/paper-gaps/.
% Every note loads this file via % Shared commands for the notes in docs/paper-gaps/.
% Every note loads this file via \input{command}. Project-wide notation
% belongs here, never in per-note redefinitions.

% --- Project configuration ------------------------------------------------
\newcommand{\projectIssueBase}{https://github.com/LionSR/QICLean/issues/}
\newcommand{\projectPrBase}{https://github.com/LionSR/QICLean/pull/}
\newcommand{\projectDocBase}{https://sirui-lu.com/QICLean/docs/}
% ---------------------------------------------------------------------------

\DeclareUnicodeCharacter{2080}{\ifmmode{}_{0}\else\textsubscript{0}\fi}
\DeclareUnicodeCharacter{2081}{\ifmmode{}_{1}\else\textsubscript{1}\fi}
\DeclareUnicodeCharacter{2082}{\ifmmode{}_{2}\else\textsubscript{2}\fi}
\DeclareUnicodeCharacter{2099}{\ifmmode{}_{n}\else\textsubscript{n}\fi}

\newcommand{\Lean}{\textsc{Lean}}

% A formal declaration name, upright sans, breakable at underscores.
\ExplSyntaxOn
\NewDocumentCommand{\leanid}{m}{
  \ifmmode
    \textsf{\detokenize{#1}}
  \else
    \group_begin:
    \urlstyle{sf}
    \tl_set:Nn \l_tmpa_tl {#1}
    \tl_replace_all:Nnn \l_tmpa_tl {\_} {_}
    \exp_args:NV \path \l_tmpa_tl
    \group_end:
  \fi
}
\ExplSyntaxOff

% Traceability links; use in footnotes, not in the running prose.
\newcommand{\ghissue}[1]{\href{\projectIssueBase#1}{GitHub issue \##1}}
\newcommand{\ghpr}[1]{\href{\projectPrBase#1}{PR \##1}}
\newcommand{\doclink}[1]{\href{\projectDocBase#1.html}{\path{#1}}}

% Domain notation, using \providecommand so a note-local refinement stays
% possible.
\providecommand{\tr}{\operatorname{tr}}

% Dirac notation and the identity operator, as in blueprint/src/macros.
\usepackage{dsfont}
\providecommand{\ket}[1]{|#1\rangle}
\providecommand{\bra}[1]{\langle #1|}
\providecommand{\braket}[2]{\langle #1 | #2 \rangle}
\providecommand{\ketbra}[2]{|#1\rangle\!\langle #2|}
\providecommand{\Id}{\mathds{1}}
\providecommand{\one}{\mathds{1}}
\providecommand{\C}{\mathbb{C}}
\providecommand{\MN}[1]{M_{#1}(\C)}
\providecommand{\MD}{M_D(\C)}

% Verdict marker: \gapnote{<kind>}{<status>}. Kind is the policy
% classification (clarification, local-correction, scope-restriction,
% unfaithful, false-source, open-gap); status is open, wip (elimination
% underway), resolved, or historical. Machine-read by the site index;
% prints a verdict line.
\newcommand{\gapnote}[2]{\par\noindent\textbf{Verdict:} #1 (#2).\par}
. Project-wide notation
% belongs here, never in per-note redefinitions.

% --- Project configuration ------------------------------------------------
\newcommand{\projectIssueBase}{https://github.com/LionSR/QICLean/issues/}
\newcommand{\projectPrBase}{https://github.com/LionSR/QICLean/pull/}
\newcommand{\projectDocBase}{https://sirui-lu.com/QICLean/docs/}
% ---------------------------------------------------------------------------

\DeclareUnicodeCharacter{2080}{\ifmmode{}_{0}\else\textsubscript{0}\fi}
\DeclareUnicodeCharacter{2081}{\ifmmode{}_{1}\else\textsubscript{1}\fi}
\DeclareUnicodeCharacter{2082}{\ifmmode{}_{2}\else\textsubscript{2}\fi}
\DeclareUnicodeCharacter{2099}{\ifmmode{}_{n}\else\textsubscript{n}\fi}

\newcommand{\Lean}{\textsc{Lean}}

% A formal declaration name, upright sans, breakable at underscores.
\ExplSyntaxOn
\NewDocumentCommand{\leanid}{m}{
  \ifmmode
    \textsf{\detokenize{#1}}
  \else
    \group_begin:
    \urlstyle{sf}
    \tl_set:Nn \l_tmpa_tl {#1}
    \tl_replace_all:Nnn \l_tmpa_tl {\_} {_}
    \exp_args:NV \path \l_tmpa_tl
    \group_end:
  \fi
}
\ExplSyntaxOff

% Traceability links; use in footnotes, not in the running prose.
\newcommand{\ghissue}[1]{\href{\projectIssueBase#1}{GitHub issue \##1}}
\newcommand{\ghpr}[1]{\href{\projectPrBase#1}{PR \##1}}
\newcommand{\doclink}[1]{\href{\projectDocBase#1.html}{\path{#1}}}

% Domain notation, using \providecommand so a note-local refinement stays
% possible.
\providecommand{\tr}{\operatorname{tr}}

% Dirac notation and the identity operator, as in blueprint/src/macros.
\usepackage{dsfont}
\providecommand{\ket}[1]{|#1\rangle}
\providecommand{\bra}[1]{\langle #1|}
\providecommand{\braket}[2]{\langle #1 | #2 \rangle}
\providecommand{\ketbra}[2]{|#1\rangle\!\langle #2|}
\providecommand{\Id}{\mathds{1}}
\providecommand{\one}{\mathds{1}}
\providecommand{\C}{\mathbb{C}}
\providecommand{\MN}[1]{M_{#1}(\C)}
\providecommand{\MD}{M_D(\C)}

% Verdict marker: \gapnote{<kind>}{<status>}. Kind is the policy
% classification (clarification, local-correction, scope-restriction,
% unfaithful, false-source, open-gap); status is open, wip (elimination
% underway), resolved, or historical. Machine-read by the site index;
% prints a verdict line.
\newcommand{\gapnote}[2]{\par\noindent\textbf{Verdict:} #1 (#2).\par}
. Project-wide notation
% belongs here, never in per-note redefinitions.

% --- Project configuration ------------------------------------------------
\newcommand{\projectIssueBase}{https://github.com/LionSR/QICLean/issues/}
\newcommand{\projectPrBase}{https://github.com/LionSR/QICLean/pull/}
\newcommand{\projectDocBase}{https://sirui-lu.com/QICLean/docs/}
% ---------------------------------------------------------------------------

\DeclareUnicodeCharacter{2080}{\ifmmode{}_{0}\else\textsubscript{0}\fi}
\DeclareUnicodeCharacter{2081}{\ifmmode{}_{1}\else\textsubscript{1}\fi}
\DeclareUnicodeCharacter{2082}{\ifmmode{}_{2}\else\textsubscript{2}\fi}
\DeclareUnicodeCharacter{2099}{\ifmmode{}_{n}\else\textsubscript{n}\fi}

\newcommand{\Lean}{\textsc{Lean}}

% A formal declaration name, upright sans, breakable at underscores.
\ExplSyntaxOn
\NewDocumentCommand{\leanid}{m}{
  \ifmmode
    \textsf{\detokenize{#1}}
  \else
    \group_begin:
    \urlstyle{sf}
    \tl_set:Nn \l_tmpa_tl {#1}
    \tl_replace_all:Nnn \l_tmpa_tl {\_} {_}
    \exp_args:NV \path \l_tmpa_tl
    \group_end:
  \fi
}
\ExplSyntaxOff

% Traceability links; use in footnotes, not in the running prose.
\newcommand{\ghissue}[1]{\href{\projectIssueBase#1}{GitHub issue \##1}}
\newcommand{\ghpr}[1]{\href{\projectPrBase#1}{PR \##1}}
\newcommand{\doclink}[1]{\href{\projectDocBase#1.html}{\path{#1}}}

% Domain notation, using \providecommand so a note-local refinement stays
% possible.
\providecommand{\tr}{\operatorname{tr}}

% Dirac notation and the identity operator, as in blueprint/src/macros.
\usepackage{dsfont}
\providecommand{\ket}[1]{|#1\rangle}
\providecommand{\bra}[1]{\langle #1|}
\providecommand{\braket}[2]{\langle #1 | #2 \rangle}
\providecommand{\ketbra}[2]{|#1\rangle\!\langle #2|}
\providecommand{\Id}{\mathds{1}}
\providecommand{\one}{\mathds{1}}
\providecommand{\C}{\mathbb{C}}
\providecommand{\MN}[1]{M_{#1}(\C)}
\providecommand{\MD}{M_D(\C)}

% Verdict marker: \gapnote{<kind>}{<status>}. Kind is the policy
% classification (clarification, local-correction, scope-restriction,
% unfaithful, false-source, open-gap); status is open, wip (elimination
% underway), resolved, or historical. Machine-read by the site index;
% prints a verdict line.
\newcommand{\gapnote}[2]{\par\noindent\textbf{Verdict:} #1 (#2).\par}

\title{The weight range and coefficient for zero-floor maximizing filters}
\author{}
\date{2026-10-16}
\begin{document}
\maketitle
\gapnote{scope-restriction}{open}

\section{Source assertion}

The companion development of Ref.~\cite{OpenAI2026PEPSFiltersEnergy}
constructs a norm-maximizing product of nested filters on a square grid.
For a local dimension $q$, interaction bound $J$, normalized unique ground
vector $\Omega$ of energy $E_0$, and full-system gap $\Delta$, it permits
weights $0<a_j\le1$. Each edge crosses at most one filtering region $D_j$.
If $b_j$ counts the edges crossing $D_j$, the selected normalized output
$\phi$ satisfies
\begin{align}
  |\langle\phi,H\phi\rangle-E_0|
    &\le q^2J\sum_j b_j a_j^2,
    \label{eq:zero_floor_source_energy}\\
  \Delta(1-|\langle\Omega,\phi\rangle|^2)
    &\le q^2J\sum_j b_j a_j^2.
    \label{eq:zero_floor_source_overlap}
\end{align}
The source first takes a compact limit of regularized maximizing families,
proves that the limiting maximal output norm is positive, and then passes
the energy and overlap bounds through normalization.\footnote{Source
\leanid{exists_maximizer_energy}, lines 62--109 of
\href{https://github.com/openai/math/blob/adc7f1241b42e322a6451854ab7e4b4c146bf78a/lean/OAI/MathematicalPhysics/PEPSFilters/EnergyOptimizer.lean}{EnergyOptimizer.lean}.}

\section{The restricted estimate}

The present result applies to arbitrary finite sites and supported
Hermitian terms $H=\sum_i h_i$, with $\|h_i\|\le c_0$ and
$H\Omega=E_0\Omega$, where $\Omega\ne0$. It assumes nested regions,
$0<a_j\le1/2$, and that each support $S_i$ is split by at most one region.
Write $d_i=\dim\mathcal H_{S_i}$ and let $C_j$ be the terms split by $D_j$.
One floor-zero norm maximizer has output $\psi\ne0$ and obeys
\begin{align}
 B&=\sum_j4a_j^2\sum_{i\in C_j}d_i^2c_0,
   \label{eq:zero_floor_coarse_constant}\\
 \operatorname{Re}\langle\psi,H\psi\rangle-E_0\|\psi\|^2
   &\le B\|\psi\|^2.
   \label{eq:zero_floor_homogeneous}
\end{align}
Consequently its unit normalized output has energy excess at most $B$.
A full-system gap yields an overlap bound for the same witness when
$\Omega$ is normalized. The energy assertion alone does not assume that
$E_0$ is the ground energy, so it states an upper bound rather than an
absolute-value estimate.\footnote{Formal declarations
\leanid{Entropy.exists_max_feasible_chain_zero_floor_energy} and
\leanid{Entropy.exists_max_feasible_chain_zero_floor_normalized_energy}
in \doclink{QICLean/Entropy/FilterZeroFloorEnergy}.}

Both the weight restriction and the coefficient in
(\ref{eq:zero_floor_coarse_constant}) enter before taking the limit:
they are the hypotheses and conclusion of the available positive-floor
maximizing-family energy estimate. The compactness argument introduces
neither restriction. It instead passes the closed inequality
(\ref{eq:zero_floor_homogeneous}) to a limiting maximizer. A fixed uniform
positive-floor competitor shows that this maximizer has nonzero output,
so division by its squared norm is legitimate.\footnote{The positive-floor
provider is \leanid{Entropy.re_inner_sub_le_of_max} in
\doclink{QICLean/Entropy/FilterChainEnergy}.}

\section{Remaining comparison}

For two-site interactions of local dimension $q$, using the full support
dimension in (\ref{eq:zero_floor_coarse_constant}) gives $4q^4$, not the
$q^2$ in (\ref{eq:zero_floor_source_energy}). Thus the zero-floor result
does not establish (\ref{eq:zero_floor_source_energy}) throughout the
source's weight range, nor its complete physical conclusion.

To remove this restriction, derive a maximizing-family positive-floor
estimate with weights at most one and the inside-support dimension from
the existing sharp chain inequality. Then use the same closed-condition
limit theorem, identify the grid terms and feasible filters with the source
representations, and apply the ground-state lower bound and full-system gap.
The geometric producer additionally needs its uniform collar, weight and
representation arguments.\footnote{Selective-port comparison remains
tracked in \href{https://github.com/LionSR/TNLean/issues/8741}{TNLean issue 8741}.}
The verdict is a scope restriction of the companion-source comparison,
not an error in its assertion. The restricted theorem is useful on its
stated domain; the source comparison remains open.

\bibliographystyle{alpha}
\bibliography{references}
\end{document}
